\documentclass[11pt]{article}
\usepackage[margin=1in]{geometry}
\usepackage{amsmath,amssymb,booktabs,hyperref,graphicx}
\title{Independent Replication --- arXiv:1808.03623\\
\large Mitigating Algorithmic Errors in Hamiltonian Simulation}
\author{Ollie (subagent) \\ Wave: QC-100}
\date{2026-07-03 (backfill 2026-07-06)}

\begin{document}
\maketitle

\section{Paper}
Endo, Zhao, Li, Benjamin, Yuan, \emph{Mitigating algorithmic errors in
Hamiltonian simulation}, arXiv:1808.03623 (2018). The paper extends
Richardson-style extrapolation --- previously used to mitigate
\emph{physical} gate errors --- to the Trotter step count $1/N$, thereby
mitigating \emph{algorithmic} (product-formula) error. The 5-qubit
transverse-field Ising benchmark of Sec.~V is used to demonstrate an
orders-of-magnitude reduction in residual Trotter error via 2- and
3-point extrapolation.

\section{Scope tested}
This replication targets the \emph{algorithmic-error-mitigation} core
(Sec.~IV plus the noise-free portion of Fig.~3). Claims C1--C4 (Trotter
series structure, linear 2-pt identity, Richardson 3-pt identity, and
quantitative Fig.-3(a) reduction) are tested. Claims C5--C6 (physical
Pauli-noise Fig.~3(b,c) and combined noise+algorithmic extrapolation)
are explicitly \emph{not} re-run in this budget.

\section{Method (independent)}
Qiskit 2.5.0 + qiskit-aer 0.17.2 on macOS. Hamiltonian built as
\texttt{SparsePauliOp}; exact reference $\langle X_1\rangle_{\text{exact}}$
computed by \texttt{scipy.linalg.expm} on the 32-dim statevector
(\emph{not} taken from the paper). First-order Trotter step assembled
from standard CNOT--$R_z$--CNOT decompositions for
$e^{-iJ\Delta t Z_iZ_{i+1}}$ and $R_x(2B\Delta t)$ for
$e^{-iB\Delta t X_i}$; $N$ steps composed then statevector-evolved.
Total wallclock $<2\,\mathrm{s}$.

\section{Headline result}
For the paper's Fig.-3(a) triple $(N_1,N_2,N_3)=(15,20,25)$:

\begin{center}\begin{tabular}{lrl}
\toprule
Method & $|\text{err}|$ & vs.\ raw $N{=}25$ \\
\midrule
No mitigation, $N{=}25$ & $4.97{\times}10^{-4}$ & $1.00\times$ \\
Linear 2-pt $(15,25)$ & $1.88{\times}10^{-3}$ & $0.26\times$ (worse) \\
\textbf{Richardson 3-pt $(15,20,25)$} & $\mathbf{2.67{\times}10^{-5}}$ & $\mathbf{18.6\times}$ \\
\bottomrule
\end{tabular}\end{center}

An independent quadratic fit
$A(1/N) = a_0 + a_1/N + a_2/N^2$ to nine data points recovers
$a_0 = 0.6729581$ (error $2.97{\times}10^{-5}$ from exact),
$a_1 \approx 0.017$, $a_2 \approx -0.72$, corroborating C1.

\section{Verdict: REPLICATED}
The technique's central noise-free claim reproduces on an
independently-built Qiskit circuit against an independent exact
reference. Fig.~3(a) qualitative and (within the paper's own step
triple) quantitative behavior matches: Richardson 3-pt yields the
promised orders-of-magnitude reduction (18.6$\times$ here); linear 2-pt
underperforms at this triple as the paper itself notes when $1/N^2$
dominates the residual.

\section{Honest critique (what is \emph{not} shown here)}
\begin{itemize}
  \item \textbf{Physical noise not simulated.} Fig.~3(b,c), which show
        the practical case of a real device where $N_{\text{opt}}$
        emerges from the trade-off between algorithmic error
        ($\downarrow$ with $N$) and physical error ($\uparrow$ with
        $N$), was not reproduced. Without this, we cannot confirm the
        paper's specific claim that $N_{\text{opt}}=25$ for its Pauli
        channel, nor that the combined algorithmic + physical
        extrapolation dominates either mitigation alone. This is the
        \emph{practically important} regime; skipping it means our
        REPLICATED verdict speaks only to the mathematical identity,
        not to end-to-end usefulness on noisy hardware.
  \item \textbf{Only one Hamiltonian.} The paper's benchmark (5-qubit
        TFIM, $J{=}3, B{=}2, t{=}0.5$) is the sole test. We did not
        probe Heisenberg, molecular (H$_2$/LiH), or non-commuting
        chemistry Hamiltonians where higher-order Trotter error
        structure is more delicate. Generalization is assumed, not
        demonstrated.
  \item \textbf{No comparison against alternative Trotter-error
        suppression.} qDRIFT, randomized compilations (Campbell 2019),
        higher-order product formulas (Suzuki-4), and multi-product
        formulas (Low--Kliuchnikov--Wiebe) all target the same
        algorithmic error. The paper (2018) predates several of these;
        our replication also does not benchmark them. A reader asking
        "should I use this vs.\ qDRIFT?" gets no answer from either
        the paper or from us.
  \item \textbf{$1/N^k$ series is assumed.} The Richardson fit relies
        on the leading Trotter expansion being polynomial in $1/N$.
        For time-dependent Hamiltonians, non-analytic error structure
        (e.g., commutator-scaling terms), or when $t$ is not small
        compared with $\|H\|^{-1}$, this assumption can break. Not
        stress-tested.
  \item \textbf{Ordering sensitivity untested.} The choice of Trotter
        ordering (e.g.\ Strang vs.\ Lie, ZZ-block first vs.\ X-block
        first) changes $a_1, a_2$ and therefore the residual after
        extrapolation. We used one canonical order.
  \item \textbf{Numerical floor at $N \gtrsim 75$.} The raw error stops
        decreasing near $\sim 5{\times}10^{-5}$ at large $N$, which we
        attribute to floating-point noise in the compiled statevector
        pipeline. This is a replication-side artifact, not a paper
        artifact, but it caps the maximum residual improvement we can
        demonstrate.
\end{itemize}

\section{Environment and evidence}
Qiskit 2.5.0, qiskit-aer 0.17.2, numpy 2.4.3, scipy 1.18.0.
Evidence: \texttt{report/evidence/\{results.json, scaling\_curve.json,
run.log\}}. Code: \texttt{code/replicate\_algo\_error\_mitigation.py},
\texttt{code/error\_scaling\_curve.py}. Full method and prose:
\texttt{report/REPORT.md}.

\end{document}
